% Pursuing Stacks (cotes 134-2 to 134-8): archive cote 134-8, site page 4 (PDF page 5), typescript page 558.
% typescript. Draft compared: ps7.tex:133-208.
% AI-transcribed by Claude Opus 5.5 (claude-opus-5-5, adaptive thinking, effort high), Message Batches API, batch ps4 (msgbatch_012Aoc1NuskATPsyBeKSYK7J), 2026-11-12; reviewed against both supplied scan halves; see CHECK_DECISIONS.tsv.
\scanpage{134-8/4}
\mspage{558}

presheaf on A,  and
\[ \text{(2)}\qquad \mathrm{LH}. = \mathrm{Lp}_{!}^{ab} : D^{-}(A^{\hand{\wedge}}_{ab}) \hand{\longrightarrow} D^{-}(Ab) \hand{,} \quad \text{inducing}\ \margin{][} D.(A^{\hand{\wedge}}_{ab}) \hand{\longrightarrow} D.(Ab)\hand{,} \]
is it's left derived functor (computed using projective resolutions of
complexes bounded above in $A^{\hand{\wedge}}_{ab}$ ),  composing (2) with the tautological
inclusion functor
\[ A^{\hand{\wedge}}_{ab} \hand{\hookrightarrow} D^{\deleted{\illegible{1}}}.(A^{\hand{\wedge}}_{ab}) \hand{\hookrightarrow} D^{-}(A^{\hand{\wedge}}_{ab}) \quad , \]
we get a canonical functor
\[ \text{(3)}\qquad \mathrm{Lp}_{!}^{ab}\big| A^{\hand{\wedge}}_{ab} : \ A^{\hand{\wedge}}_{ab} \hand{\longrightarrow} D.(Ab) \ , \]
which may be expressed,  in terms of an integrator  $L.^{B}$  for  $B = A^{o}$\hand{,}
i.e. a projective resolution of the constant presheaf $\underline{Z}_B$  in $B^{\hand{\wedge}}_{ab}$ ,  as
the composition
\[ \text{(4)}\qquad A^{\hand{\wedge}}_{ab} \hand{\longrightarrow} Ch.(Ab) \hand{\longrightarrow} D.(Ab) \quad , \]
where the first arrow is
\[ \text{(5)}\qquad F \hand{\longmapsto} F *_{\underline{Z}} L.^{B} \ : \deleted{\illegible{1}}\ A^{\hand{\wedge}}_{ab} \hand{\longrightarrow} Ch.(Ab) \qquad , \]
and the second is the canonical localization functor.   The functors
(2) and (3) may be viewed as the (total) homology functors of A,  with
coe\over{\unsure{ffi}}{ffi}cients in \xout{Objects of}  complexes of abelian presheaves,  resp.
in abelian presheaves simply.  When we focus attention on the latter,
we may introduce in $A^{\hand{\wedge}}_{ab}$ the set of arrows which \xout{are transformed} become
isomorphisms under the \xout{functor} total homology functor (3),  let's call
them ``\uline{abelian} \uline{weak} \uline{equivalences}'' in $A^{\hand{\wedge}}_{ab}$ ,  not to be confused with the
notion of quasi-isomorphism for a map \over[typed]{\unsure{p}}{b}etween complexes in $A^{\hand{\wedge}}_{ab}$ .  Let's
denote by
\[ \text{(6)}\qquad \mathrm{Hotab}_A \ = (\over{\unsure{W}}{W}_A^{ab})^{-1} A^{\hand{\wedge}}_{ab} \]
the corresponding localized category of $A^{\hand{\wedge}}_{ab}$ ,  where  $W_A^{ab}$  denotes the
set of abelian weak equivalences in $A^{\hand{\wedge}}_{ab}$ .   Thus the choice of an inte-%
grator  $L.^{B\deleted{\illegible{1}}}$  for A (i.e. a cointegrator for B) gives rise to a commutative
diagram of functors
\diagram{
% Hand-drawn arrows are marked blue; typed objects remain plain.
\text{(7)}\qquad
\begin{tikzcd}
A^{\hand{\wedge}}_{ab} \arrow[color=blue,r] \arrow[color=blue,d] & Ch.(Ab) \arrow[color=blue,d] \\
\mathrm{Hotab}_A \arrow[color=blue,r] & D.(Ab) \overset{def}{=} \mathrm{Hotab}
\end{tikzcd}
\qquad ,
}
where the \xout{second} lower horizontal arrow is defined via (3) \hand{(}independently
of the choice of $L.^{B}$\hand{)},  the vertical arrows being the localization func-%
tors.   Beware that even when  $A=\hand{\Delta}$  and $L.^{B}$  is the usual, ``standar\over{t}{d}''
\margin{\textcircled{$L.^B$}}
\margin{4}

% Visual collation: both supplied image halves inspected; see CHECK_DECISIONS_134_7_8.tsv.
